\documentclass[conference]{IEEEtran}

\usepackage[utf8]{inputenc}
\usepackage[T1]{fontenc}
\usepackage{newtxtext,newtxmath}
\usepackage{amsmath,bm}
\usepackage{mathtools}
\usepackage{booktabs}
\usepackage{siunitx}
\usepackage{cite}
\usepackage{listings}
\usepackage{xcolor}
\usepackage[breaklinks=true,hidelinks]{hyperref}

% ---- Code listing style -----------------------------------------------
\lstdefinestyle{cppstyle}{
    language=C++,
    basicstyle=\ttfamily\scriptsize,
    keywordstyle=\color{blue!70!black}\bfseries,
    commentstyle=\color{gray}\itshape,
    stringstyle=\color{orange!80!black},
    breaklines=true,
    breakatwhitespace=true,
    frame=single,
    rulecolor=\color{black!30},
    xleftmargin=4pt,
    xrightmargin=4pt,
    numbers=none,
    tabsize=2,
}
\lstset{style=cppstyle}

% -----------------------------------------------------------------------

\title{Embedded-Friendly IMU Calibration:\\
Ellipsoid Fitting, Temperature Compensation,\\
and Robust Estimation in C\texttt{++}/Eigen}

\author{%
    \IEEEauthorblockN{Mikhail S.\ Grushinskiy}
    \IEEEauthorblockA{Independent Researcher\\
    Bareboat Necessities Project\\
    \texttt{mgrouch@users.github.com}}
}

\begin{document}
    \maketitle

% -----------------------------------------------------------------------
    \begin{abstract}
        This note documents \texttt{imu\_cal}, a header-only C\texttt{++}
        library for in-situ calibration of the three sensors of an Inertial
        Measurement Unit (IMU): accelerometer, magnetometer, and gyroscope,
        for Arduino-class targets using the Eigen linear-algebra library.
        The magnetometer and the library accelerometer calibrator fit an
        ellipsoid with a robust algebraic 10-parameter least-squares solver,
        trimmed iterative reweighting and positive-definite projection.
        The AtomS3R wizard calibrates the accelerometer with a guided manual
        procedure: ten static holds (six faces and four tilted corners chosen
        from an information study), bounded targeted extra holds, and three
        later rechecks. A geometric, double-precision Levenberg--Marquardt fit
        keeps the full symmetric positive-definite scale/cross-axis matrix,
        estimates the bias from measured vectors and the gravity magnitude
        only, and fits a linear thermal bias slope jointly when the data
        determine it after the static parameters. Every result is checked on
        held-out holds and re-validated as the stored float coefficients.
        A deterministic simulation campaign compares the deployed procedure,
        the new capture with the old fitter, and the new procedure against
        known truth.
    \end{abstract}

    \begin{IEEEkeywords}
        IMU calibration, ellipsoid fitting, robust estimation,
        temperature compensation, Eigen, embedded systems, Arduino,
        accelerometer, magnetometer, gyroscope
    \end{IEEEkeywords}

% -----------------------------------------------------------------------
    \section{Introduction}

    Low-cost MEMS IMUs suffer from well-known systematic errors:
    scale-factor mismatches, inter-axis cross-sensitivity, constant bias,
    and temperature-dependent bias drift \cite{Frosio2009}.
    Without calibration, these errors accumulate rapidly in attitude
    and heading estimates.

    Classical batch calibration techniques require a precision turntable
    or a known-orientation jig.
    In contrast, \emph{in-situ} or \emph{motional} calibration methods
    collect raw sensor data while the device is rotated freely and
    then recover the error model from the geometry of the measurements
    \cite{Vasconcelos2011}.
    For the accelerometer and magnetometer, the underlying observation
    is that the true sensor output should lie on a sphere of known
    radius ($g$ for acceleration, the local field magnitude $B$ for
    the magnetometer).
    Systematic errors deform this sphere into an ellipsoid; recovering
    the ellipsoid and inverting the distortion yields the calibration map.

    \texttt{imu\_cal} implements this approach for all three sensor axes
    in dependency-light C\texttt{++} headers suitable
    for deployment on Arduino-class hardware.
    For the accelerometer the manual AtomS3R procedure of
    Section~\ref{sec:manual} replaces free rotation by guided static holds,
    because gravity is only a reference while the device is at rest.
    The library uses Eigen \cite{Eigen} for all matrix operations
    and exposes a minimal, typed API.

% -----------------------------------------------------------------------
    \section{Sensor Error Models}

    \subsection{Accelerometer and Magnetometer}

    Let $\mathbf{v}_\text{raw} \in \mathbb{R}^3$ be the raw sensor reading.
    The standard affine error model is
    \begin{equation}
        \mathbf{v}_\text{raw} = \mathbf{M}\,\mathbf{v}_\text{true} + \mathbf{b}
        + \boldsymbol{\eta},
    \end{equation}
    where $\mathbf{M} \in \mathbb{R}^{3\times 3}$ captures scale factors and
    cross-sensitivity, $\mathbf{b} \in \mathbb{R}^3$ is the constant bias, and
    $\boldsymbol{\eta}$ is noise.
    The calibrated output is therefore
    \begin{equation}
        \mathbf{v}_\text{cal} = \mathbf{S}\bigl(\mathbf{v}_\text{raw} -
        \mathbf{b}(T)\bigr),
        \label{eq:cal_model}
    \end{equation}
    where $\mathbf{S} = \mathbf{M}^{-1}$ is the whitening matrix and
    $\mathbf{b}(T)$ is the temperature-dependent bias (Section~\ref{sec:tempbias}).

    \subsection{Gyroscope}

    The gyroscope error model used here is the simpler bias-only form,
    because scale-factor and cross-sensitivity calibration requires a
    precision rate rig:
    \begin{equation}
        \boldsymbol{\omega}_\text{cal} = \mathbf{S}\bigl(\boldsymbol{\omega}_\text{raw}
        - \mathbf{b}(T)\bigr),
        \quad \mathbf{S} = \mathbf{I}_3.
    \end{equation}
    The identity scale matrix is retained in the struct so that a measured-rate
    rig can upgrade it without changing the API.

    \subsection{Temperature Bias Model}
    \label{sec:tempbias}

    Gyroscope and accelerometer biases are strongly temperature-dependent.
    The library models this as a first-order linear expansion about a
    reference temperature $T_0$ (\SI{25}{\degreeCelsius} by default):
    \begin{equation}
        \mathbf{b}(T) = \mathbf{b}_0 + \mathbf{k}(T - T_0),
        \label{eq:tempbias}
    \end{equation}
    with intercept $\mathbf{b}_0 \in \mathbb{R}^3$ and slope
    $\mathbf{k} \in \mathbb{R}^3$. The gyroscope and the library
    accelerometer calibrator estimate them from per-temperature-bin centers;
    the manual accelerometer procedure fits them through the measurement
    model (Section~\ref{sec:thermal}).
    At runtime $T$ is clamped to $[T_\text{lo}, T_\text{hi}]$ (unbounded by
    default) and a missing (non-finite) temperature evaluates
    $\mathbf{b}_0$.

% -----------------------------------------------------------------------
    \section{Ellipsoid Fitting}
    \label{sec:ellipsoid}

    \subsection{Algebraic Form}

    A general ellipsoid centred at $\mathbf{b}$ satisfying
    $\|\mathbf{A}(\mathbf{x}-\mathbf{b})\| = R$ can be written in the
    quadratic algebraic form
    \begin{equation}
        \mathbf{x}^\top \mathbf{Q}\,\mathbf{x}
        + 2\mathbf{q}^\top\mathbf{x} + c = 0,
    \end{equation}
    where $\mathbf{Q} \in \mathbb{R}^{3\times 3}$ is symmetric positive
    definite, $\mathbf{q} \in \mathbb{R}^3$, and $c \in \mathbb{R}$.
    Expanding the defining equation:
    \begin{align}
        \mathbf{Q} &= \mathbf{A}^\top\mathbf{A}, \\
        \mathbf{q} &= -\mathbf{Q}\mathbf{b}, \\
        c &= \mathbf{b}^\top\mathbf{Q}\mathbf{b} - R^2.
    \end{align}

    \subsection{10-Parameter Regression}
    \label{sec:regression}

    Define the feature vector for sample $\mathbf{x}_i = (X,Y,Z)^\top$:
    \begin{equation}
        \mathbf{d}_i = \bigl[X^2,\,Y^2,\,Z^2,\,2XY,\,2XZ,\,2YZ,\,
        2X,\,2Y,\,2Z,\,1\bigr]^\top.
    \end{equation}
    The algebraic constraint becomes $\mathbf{d}_i^\top\mathbf{p} = 0$,
    where $\mathbf{p} \in \mathbb{R}^{10}$ encodes $\mathbf{Q}$, $\mathbf{q}$, $c$.
    Because the constraint is homogeneous, we fix the normalisation by
    setting the right-hand side to 1 (corresponding to $c = 1$), leading
    to the linear system
    \begin{equation}
        \mathbf{H}\,\mathbf{p} = \mathbf{g},
        \quad
        \mathbf{H} = \sum_i \mathbf{d}_i\mathbf{d}_i^\top,
        \quad
        \mathbf{g} = \sum_i \mathbf{d}_i.
        \label{eq:lsq}
    \end{equation}
    A small Tikhonov ridge
    $\lambda = \lambda_\text{rel}\!\left(\tfrac{\mathrm{tr}(\mathbf{H})}{10}
    + 10^{-12}\right)$
    is added to the diagonal of $\mathbf{H}$ before solving via LDLT
    decomposition:
    \begin{equation}
    (\mathbf{H} + \lambda\mathbf{I})\,\mathbf{p} = \mathbf{g}.
    \label{eq:ridge}
    \end{equation}
    The default ridge scale is $\lambda_\text{rel} = 10^{-6}$, which is
    harmless for normal data but prevents numerical failure with
    near-planar samples.

    \subsection{Model Reconstruction}

    From the solution $\mathbf{p}$, the centre and shape matrix are recovered as:
    \begin{align}
        \mathbf{b} &= -\mathbf{Q}^{-1}\mathbf{q}, \\
        s &= 1 + \mathbf{b}^\top\mathbf{Q}\mathbf{b} - c, \\
        \mathbf{M} &= \mathbf{Q}/s.
    \end{align}
    The matrix $\mathbf{M}$ is projected to symmetric positive definite (SPD)
    via eigenvalue clamping (Section~\ref{sec:spd}), then Cholesky-factored
    $\mathbf{M} = \mathbf{U}^\top\mathbf{U}$ to obtain the unit-sphere
    whitening matrix $\mathbf{A}_\text{unit} = \mathbf{U}$.
    The final calibration matrix scaled to the target radius $R$ is
    \begin{equation}
        \mathbf{A} = \mathbf{A}_\text{unit}\,R.
    \end{equation}

% -----------------------------------------------------------------------
    \section{Robust Fitting}
    \label{sec:robust}

    A single pass of algebraic fitting is sensitive to outliers, ambient
    interference near the magnetometer, and the unavoidable non-uniformity
    of free-motion data collection.
    The library therefore wraps the linear solve inside an iterative
    trimming loop.

    \subsection{Algorithm Overview}

    \begin{enumerate}
        \item Initialise all sample weights as inliers.
        \item Solve \eqref{eq:lsq}--\eqref{eq:ridge} on the current inlier set.
        \item Compute residuals $e_i = \|\mathbf{A}(\mathbf{x}_i - \mathbf{b})\| - R$.
        \item Compute the Median Absolute Deviation (MAD) of $\{e_i\}$.
        \item Mark as outliers any sample with $|e_i| > \tau$, where
        $\tau = \min(\tau_\text{MAD},\,\tau_\text{quantile})$.
        \item Repeat from step 2 for a fixed number of iterations
        (default 3).
        \item Final solve on the converged inlier set.
    \end{enumerate}

    \subsection{MAD Threshold}

    The MAD-based threshold is
    \begin{equation}
        \tau_\text{MAD} = 5.2\,\max(\mathrm{MAD},\,10^{-4}) + 10^{-6}.
    \end{equation}
    The factor 5.2 is chosen empirically to be permissive enough for
    real sensor noise while still rejecting gross outliers.

    \subsection{Quantile Threshold}

    A trimming fraction $\alpha$ (default \SI{15}{\percent}) is applied by
    sorting the absolute residuals and setting
    \begin{equation}
        \tau_\text{quantile} = |e|_{(\lceil(1-\alpha)n\rceil)},
    \end{equation}
    the $(1-\alpha)$-quantile of the absolute residual distribution.
    Samples above this threshold are excluded.

    \subsection{Fallback}

    If the joint threshold would leave fewer than 10 inliers, the code
    relaxes to the looser of the two thresholds.
    If that still leaves fewer than 10, fitting fails with
    \texttt{INLIERS\_TOO\_FEW}.

% -----------------------------------------------------------------------
    \section{SPD Projection}
    \label{sec:spd}

    Numerical noise can cause the recovered matrix $\mathbf{M}$ to have
    small negative eigenvalues, which would break the subsequent Cholesky
    factorisation.
    The library projects any symmetric matrix to the nearest SPD matrix
    by clamping its eigenvalues:
    \begin{equation}
        \mathbf{M} \leftarrow \mathbf{V}\,
        \mathrm{diag}\!\left(\max(|\lambda_i|,\,\epsilon)\right)\,
        \mathbf{V}^\top,
    \end{equation}
    where $\mathbf{V}$ are the eigenvectors of $\mathbf{M}$ and
    \begin{equation}
        \epsilon = \max\!\bigl(\epsilon_\text{rel}\,\lambda_\text{max},\;
        \epsilon_\text{abs}\bigr),
        \quad
        \epsilon_\text{rel} = 10^{-5},\;
        \epsilon_\text{abs} = 10^{-7}.
    \end{equation}
    This is the nearest-SPD projection under the Frobenius norm for
    symmetric matrices \cite{Higham1988}.

% -----------------------------------------------------------------------
    \section{Coverage and Degeneracy Checks}
    \label{sec:coverage}

    Before fitting, the library verifies that the collected samples
    adequately cover the sphere.
    Let $\hat{\mathbf{x}}_i = (\mathbf{x}_i - \boldsymbol{\mu})/
    \|\mathbf{x}_i - \boldsymbol{\mu}\|$ be the unit-direction vectors
    after mean-centering.

    Three conditions must hold simultaneously:
    \begin{enumerate}
        \item \textbf{Axis span}: the centred sample range on each axis
        must satisfy $\Delta_j \ge 0.90\,R_\text{expected}$.
        \item \textbf{Sign balance}: at least \SI{8}{\percent} of samples
        must fall on each side of zero on every axis.
        \item \textbf{Directional spread}: the direction covariance
        $\mathbf{C} = \frac{1}{m}\sum_i \hat{\mathbf{x}}_i
        \hat{\mathbf{x}}_i^\top$
        must satisfy $\det(\mathbf{C}) \ge 10^{-3}$.
        For a uniform sphere $\mathbf{C} \approx \tfrac{1}{3}\mathbf{I}$
        with $\det \approx 0.037$; the threshold rejects planar or
        linear sample distributions.
    \end{enumerate}
    Failure in any condition produces \texttt{DEGENERATE\_COVERAGE}.

% -----------------------------------------------------------------------
    \section{Library Accelerometer Calibrator}
    \label{sec:accel_fit}

    \texttt{AccelCalibrator} fits the ellipsoid of
    Section~\ref{sec:ellipsoid} to accelerometer samples.
    The AtomS3R wizard no longer uses it for the accelerometer
    (Section~\ref{sec:manual}); it remains available in the library and
    serves as the old fitter in the validation campaign.

    \subsection{Static Sample Gating}

    Only samples collected while the IMU is near-stationary are useful for
    accelerometer calibration.
    The \texttt{addSample} method enforces two gates:
    \begin{align}
        \bigl|\|\mathbf{a}_\text{raw}\| - g\bigr| &< \delta_a
        \quad (\delta_a = \SI{0.25}{\meter\per\second\squared}), \\
        \|\boldsymbol{\omega}_\text{raw}\| &< \delta_\omega
        \quad (\delta_\omega = \SI{0.07}{\radian\per\second}).
    \end{align}
    Samples not passing both gates are silently discarded.

    \subsection{Temperature Binning}

    Samples are partitioned into $K = 8$ temperature bins of equal width
    spanning $[T_\text{min}, T_\text{max}]$.
    Each bin with at least 30 samples is independently fitted by
    \texttt{ellipsoid\_to\_sphere\_robust} with target radius $g$.

    \subsection{Axis-Rotation Prevention}
    \label{sec:polar}

    The raw whitening matrix $\mathbf{A}$ returned by the ellipsoid fit
    may contain a rigid rotation component that would misalign the calibrated
    axes.
    To remove it, the library extracts the \emph{polar SPD factor}:
    \begin{equation}
        \mathbf{P} = \sqrt{\mathbf{A}^\top\mathbf{A}},
        \quad\text{so that}\quad
        \mathbf{A} = \mathbf{R}\,\mathbf{P},\;
        \mathbf{P} \succ 0.
    \end{equation}
    $\mathbf{P}$ captures scale and cross-sensitivity while discarding any
    rotation $\mathbf{R}$.
    Optionally, only the diagonal of $\mathbf{P}$ is retained
    (\texttt{AccelSMode::DiagonalOnly}, the default), which further
    prevents axis permutation at the cost of ignoring minor cross-axis
    terms.

    \subsection{Plausibility Gates}

    A bin's $\mathbf{S}$ matrix is accepted only if:
    \begin{itemize}
        \item all diagonal entries lie in $[0.80,\,1.25]$,
        \item the SPD condition number satisfies $\kappa(\mathbf{S}) \le 3.5$,
        \item in \texttt{PolarSPD} mode, the off-diagonal RMS satisfies
        $\text{rms}_\text{off}(\mathbf{S}) \le 0.05$.
    \end{itemize}
    These thresholds reject numerically or physically implausible fits
    (\texttt{ACCEL\_S\_UNPHYSICAL}).
    A bin contributes a bias center only after its matrix passed every
    gate, and a failed fit leaves \texttt{ok} false; an all-rejected fit
    therefore fails rather than returning the identity default.

    \subsection{Gravity Renormalisation}

    After fitting, the scalar magnitude of the calibrated output is
    matched to $g$ using a robust median of quiet-period norms:
    \begin{equation}
        \mathbf{S} \leftarrow \frac{g}{\mathrm{med}\bigl(
        \|\mathbf{S}(\mathbf{a}_i - \mathbf{b}(T_i))\|\bigr)}\,\mathbf{S}.
    \end{equation}
    The scale factor is clamped to $[0.97,\,1.03]$ (relaxing to
    $[0.90,\,1.10]$ on fallback).
    This step corrects for any remaining uniform scale error without
    rotating the axes.

% -----------------------------------------------------------------------
    \section{Manual Accelerometer Calibration}
    \label{sec:manual}

    The AtomS3R wizard (\texttt{AtomS3R\_ImuCalWizard.h}) runs the
    host-tested procedure in \texttt{imu\_calibrate/AccelCalCapture.h} and
    the fitter in \texttt{imu\_calibrate/AccelCalFit.h}.
    The procedure is manual: angles and opposite orientations are only
    approximate, modest hand motion is expected, and no fixture, rate table,
    heating or thermal waiting stage is used.

    \subsection{Model}

    \begin{equation}
        \mathbf{a}_\text{cal} = \mathbf{S}\bigl(\mathbf{a}_\text{raw}
        - \mathbf{b}_\text{ref} - \mathbf{k}\,(T - T_\text{ref})\bigr),
        \quad \mathbf{S} = \mathbf{S}^\top \succ 0.
        \label{eq:accel_full}
    \end{equation}
    $\mathbf{S}$ keeps its six independent entries: three scales and three
    symmetric cross-axis terms. A gravity-magnitude fit determines the
    inverse gain matrix only up to a rotation; the symmetric solution is its
    polar-SPD factor, the same convention as \texttt{PolarSPD}
    (Section~\ref{sec:polar}), so no frame rotation is introduced.
    At a fixed temperature nine parameters are estimated
    ($\mathbf{S}$, $\mathbf{b}_\text{ref}$); the three slopes $\mathbf{k}$
    are added only when the data determine them (Section~\ref{sec:thermal}).

    \subsection{Gravity Convention}
    \label{sec:gravity}

    Two gravity values are used, and they are kept apart.

    \emph{Standard gravity} $g_\text{std} = 9.80665$\,m/s$^2$ is the exact
    conventional value of one standard $g$. M5Unified reports the BMI270
    acceleration in nominal $g$; multiplying by $g_\text{std}$ converts that
    unit to m/s$^2$. The chain is: BMI270 counts, M5Unified acceleration in
    nominal $g$, $\times g_\text{std}$, SI acceleration, axis mapping,
    then Eq.~\eqref{eq:accel_full}. This is a unit conversion only. It does
    not mean that the static gravity where the device is calibrated is
    9.80665\,m/s$^2$ (\texttt{ImuCalCfg::g\_std}).

    \emph{Calibration gravity} $g_\text{cal}$ is the physical gravity assumed
    at the calibration site (\texttt{ImuCalCfg::g\_cal\_local}). The fit
    enforces, for every static hold,
    \begin{equation}
        \bigl\|\mathbf{S}\bigl(\mathbf{a}_\text{SI} - \mathbf{b}_\text{ref}
        - \mathbf{k}(T - T_\text{ref})\bigr)\bigr\| = g_\text{cal}.
        \label{eq:gravity_constraint}
    \end{equation}
    Its default is WGS84 normal gravity (NGA.STND.0036, formerly TR8350.2)
    at the project's default calibration site, Fair Lawn, New Jersey 07410:
    geodetic latitude $\varphi = 40.935833^\circ$\,N, longitude
    $74.117504^\circ$\,W, orthometric elevation about $+25$\,m above mean
    sea level, WGS84 ellipsoidal height about $h = -9$\,m. The Somigliana
    closed form gives normal gravity on the ellipsoid,
    \begin{equation}
        \gamma_0(\varphi) = \gamma_e\,
        \frac{1 + k\sin^2\varphi}{\sqrt{1 - e^2\sin^2\varphi}},
        \qquad k = \frac{b\,\gamma_p}{a\,\gamma_e} - 1,
    \end{equation}
    and the second-order height correction (TR8350.2 eq.~4-3) moves it to
    height $h$ above the ellipsoid:
    \begin{equation}
        \gamma(\varphi, h) = \gamma_0\Bigl[1 - \frac{2}{a}
        \bigl(1 + f + m - 2f\sin^2\varphi\bigr)h + \frac{3}{a^2}h^2\Bigr],
    \end{equation}
    with $a = 6378137$\,m, $f = 1/298.257223563$, $e^2 = 6.69437999\times
    10^{-3}$, $\gamma_e = 9.7803253359$, $\gamma_p = 9.8321849378$\,m/s$^2$,
    $m = \omega^2a^2b/GM = 0.00344978650684$. Longitude does not enter:
    normal gravity depends only on latitude and height, and no local
    gravity-field model is used. The height correction is defined in
    ellipsoidal height, so $h = -9$\,m is used, not the $+25$\,m
    orthometric elevation (which would give $1.05\times10^{-4}$\,m/s$^2$
    less). Hence $\gamma_0 = 9.8025327$ and
    \[
        g_\text{cal} = \gamma(40.935833^\circ, -9\,\text{m})
        = 9.8025605\ \text{m/s}^2,
    \]
    $4.09$\,mm/s$^2$ below $g_\text{std}$ (the GRS80 formula agrees to
    $1.5\times10^{-6}$). This is a model value, not a measurement. The
    formula is exact to $10^{-7}$\,m/s$^2$ for its inputs; $\pm5$\,m of
    height and $\pm0.001^\circ$ of latitude move it by $1.5\times10^{-5}$
    and $9\times10^{-7}$\,m/s$^2$. The real gravity differs from normal
    gravity by the local gravity disturbance, typically of order
    $10^{-4}$\,m/s$^2$, which is not modelled. Another site is configured
    with \texttt{-DATOMS3R\_CALIBRATION\_GRAVITY\_MPS2=<value>}.

    Only the gravity magnitude constrains the scale. If the fit assumes
    $g_a$ while the site has $g_t$, the fitted matrix is
    $(g_a/g_t)\,\mathbf{S}_\text{true}$ and the bias is unchanged. Forcing
    Fair Lawn gravity onto $g_\text{std}$ would therefore apply the common
    scale factor $g_\text{std}/g_\text{cal} = 1.000417$ to every calibrated
    acceleration, dynamic ones included. The estimators remove gravity from
    the calibrated specific force, $\mathbf{a}_w = \mathbf{R}\mathbf{a}_\text{cal}
    + g_e\hat{\mathbf{z}}$ (NED). For a still device
    $\|\mathbf{a}_\text{cal}\| = g_\text{cal}$, so the vertical residual is
    $g_e - g_\text{cal}$. The AtomS3R sketches therefore configure every
    estimator (OU-II, OU-III, TFG, PII, NLO, their startup proxies and
    still-detection gates) with the same $g_\text{cal}$: SI scale is kept
    and an ideal still sensor has zero translational acceleration. The
    library defaults stay at 9.80665\,m/s$^2$; the simulations use it as
    their world gravity.

    \subsection{Pose Design}

    Linearized at $\mathbf{S}=\mathbf{I}$, $\mathbf{b}=\mathbf{0}$, a hold
    whose measured up-direction is $\mathbf{u}$ observes
    $g\,\mathbf{u}^\top\delta\mathbf{S}\,\mathbf{u} - \mathbf{u}^\top\delta\mathbf{b}$.
    The six faces ($\mathbf{u}=\pm\mathbf{e}_i$) have no cross-axis
    regressors, so with faces alone the cross terms are determined only by
    in-hold wobble. Four extra directions were chosen by exhaustive search
    over the 12 edge and 8 corner directions, ranking the worst bias and
    cross-term standard deviations of the information matrix under ideal
    geometry and under random placement errors ($5^\circ$ on faces,
    $15^\circ$ on tilts). Corner sets constrain the cross terms best; the
    four screen-up corners are within a few percent of the best corner set
    and keep the screen readable during capture, so they were chosen
    (Table~\ref{tab:poses}).

    The test suite recomputes the information matrix of the shipped design
    with the fitter itself: per unit hold error the worst bias standard
    deviation is $0.71$ and the worst cross-term deviation $0.97$
    (ideal geometry), and the 95th percentiles under manual placement
    errors are $0.85$ and $1.43$. Six faces plus repeated faces are
    reported as undetermined (\texttt{INFO\_CROSS}).

    \begin{table*}[t]
        \centering
        \caption{Accelerometer holds. Up is the measured specific-force
        direction at rest in the device frame (top = end opposite the USB
        connector, right = right edge, out = out of the screen); documented
        body mapping: top $=+x$, right $=+y$, out $=-z$. Rotation is the
        display rotation relative to the reading rotation (text up on the
        highest edge).}
        \label{tab:poses}
        \renewcommand{\arraystretch}{1.1}
        \begin{tabular}{@{}rllll@{}}
            \toprule
            \# & Label & Instructions & Up (top, right, out) & Rotation \\
            \midrule
            1 & SCREEN UP & Screen faces up / Flat on the table & $(0,0,1)$ & 0 \\
            2 & SCREEN DOWN & Screen faces table / Flat on the table & $(0,0,-1)$ & 0 \\
            3 & USB UP & USB points up / Screen vertical & $(-1,0,0)$ & 2 \\
            4 & USB DOWN & USB points down / Screen vertical & $(1,0,0)$ & 0 \\
            5 & LEFT DOWN & Left edge down / Screen vertical & $(0,1,0)$ & 1 \\
            6 & RIGHT DOWN & Right edge down / Screen vertical & $(0,-1,0)$ & 3 \\
            7 & USB DN+LEFT DN & Screen up, $\sim$45 deg / USB end DOWN / Left edge DOWN & $(1,1,1)/\sqrt3$ & 0 \\
            8 & USB DN+RIGHT DN & Screen up, $\sim$45 deg / USB end DOWN / Right edge DOWN & $(1,-1,1)/\sqrt3$ & 0 \\
            9 & USB UP+LEFT DN & Screen up, $\sim$45 deg / USB end UP / Left edge DOWN & $(-1,1,1)/\sqrt3$ & 2 \\
            10 & USB UP+RIGHT DN & Screen up, $\sim$45 deg / USB end UP / Right edge DOWN & $(-1,-1,1)/\sqrt3$ & 2 \\
            \midrule
            C1--C3 & SCREEN UP, USB UP, LEFT DOWN & later rechecks (as poses 1, 3, 5) & & \\
            \bottomrule
        \end{tabular}
    \end{table*}

    Pose regions are broad: $35^\circ$ cones. Only the screen-up direction
    is taken from the documented mapping. Later regions use directions
    measured in the same session: screen down opposite pose 1, USB up
    screen-vertical and edge-aligned, left down along
    $\mathbf{u}_1\times\mathbf{u}_3$, and the corners and rechecks from the
    measured frame. An error in the assumed $x/y$ edge mapping therefore
    cannot reject correct placements. The measured frame is logged against
    the documented mapping for hardware confirmation.

    \subsection{Screen Sequence}

    \begin{enumerate}
        \item \texttt{IMU CAL}: with a saved calibration that includes the
        gyroscope, ``Tap: full calib / Tap x2: accel only / Tap x3: cancel'';
        otherwise ``Tap to begin''.
        \item For each hold, the preparation screen (reading rotation):
        title \texttt{ACCEL} (\texttt{EXTRA i/2} or \texttt{CHECK i/3}),
        the label (e.g.\ \texttt{7/10 USB DN+LEFT DN}), the full
        instructions, an optional reason (\texttt{Weak bias: USB axis},
        \texttt{Redo: inconsistent}), then ``Tap then place / Tap BtnA''.
        \item Capture screen in the pose rotation: title, label,
        instructions retained, a hint (\texttt{Place now}, \texttt{Hold
        still...}, \texttt{Hold still}, \texttt{Moving - hold still},
        \texttt{Check the pose}) and a bar that shows the 6.5\,s placement
        guidance until the device settles, then capture progress.
        \item \texttt{OK} ``label / Captured''; after a failed hold,
        \texttt{RETRY?} ``label / reason / Hold still, same pose / Tap:
        retry / Tap x3: abort'' (at most two retries per hold).
        \item \texttt{FIT} ``ACCEL / Working...'' (preliminary fit), then
        at most two \texttt{EXTRA} holds.
        \item Full wizard only: the existing \texttt{GYRO} and \texttt{MAG}
        stages.
        \item \texttt{CHECK 1/3}--\texttt{3/3}: SCREEN UP, USB UP, LEFT DOWN.
        \item \texttt{FIT} (final), \texttt{SAVE} ``Writing...'', then
        \texttt{DONE} ``A:1 G:1 M:1 / Saved OK / Bias sd 0.0025 m/s2 /
        Temp slope: learned|kept|none / Tap BtnA''.
    \end{enumerate}
    Any failure shows \texttt{FAILED} ``ACCEL / two reason lines /
    Previous cal kept''; nothing is written. All strings fit 21 columns (10
    for titles) of the $128\times128$ display at every rotation, which the
    tests check.

    \subsection{Hold Qualification}

    Samples carry \texttt{micros()} timestamps; exact repeats of the
    accelerometer and gyroscope words (stale data), non-increasing times and
    samples with $\bigl|\|\mathbf{a}\|-g\bigr| > 2\,$m/s$^2$ are skipped (a
    broad gate that keeps the bias being calibrated). Each 0.25\,s block is
    retained as its own gravity observation if it has at least five
    samples covering 60\,\% of the block without a 60\,ms gap, an RMS
    deviation below 0.25\,m/s$^2$, a per-axis range below 1.6\,m/s$^2$
    (shocks), a mean angular rate within 0.15\,rad/s of the stationary gyro
    level (the saved gyroscope bias, the gyro stage result, or the median of
    completed holds; 0.20\,rad/s before one is known), a within-block rate
    deviation below 0.30\,rad/s, and an up-direction inside the pose region.
    A 0.25\,s block smeared by 0.15\,rad/s shrinks its mean norm by about
    0.5\,mg, so these gates bound physical error without demanding
    stillness. The 6.5\,s placement interval is guidance, not a wait: after
    1\,s, three consecutive qualified blocks whose means differ by less than
    0.30\,m/s$^2$ settle the hold, and they are retained immediately.
    While capturing, a block whose mean jumps by more than 0.5\,m/s$^2$ or
    whose norm jumps by more than 0.12\,m/s$^2$ from the last retained block
    is rejected locally; four rejections in a row require settling again.
    A hold ends after 5\,s of retained blocks (rechecks: 5\,s fit plus
    2\,s verification-only), fails after 30\,s, and a sensor that delivers
    no fresh sample for 12\,s fails it at once. Blocks of one hold are never
    merged into one vector, so a slowly drifting hold keeps its spread.

    \subsection{Fit}

    Vectors are normalized by $g_\text{cal}$ (Eq.~\eqref{eq:gravity_constraint})
    and temperatures by $10\,^\circ$C;
    $\mathbf{S} = \mathbf{L}\mathbf{L}^\top$ with a log-diagonal lower
    Cholesky factor, so every iterate is SPD. The residual of a block is
    $\|\mathbf{S}(\mathbf{x}_i - \boldsymbol\beta - \boldsymbol\kappa\tau_i)\| - 1$.
    Each hold carries total weight one, split equally over its fit blocks,
    so correlated blocks and longer holds do not count as extra geometry.
    Levenberg--Marquardt runs in double precision inside six IRLS passes
    with Huber weights (knee 2 robust sigmas, MAD scale floored at
    2\,mm/s$^2$); a block beyond $\max(8\sigma, 0.08\,\text{m/s}^2)$ gets zero
    weight, but every hold keeps its best block. The start is
    $\mathbf{S} = s\mathbf{I}$, $\mathbf{b}=\mathbf{0}$ with $s$ from the
    mean norm; no simpler model replaces the full-matrix result.

    The information matrix uses the physical parameters (six entries of
    $\mathbf{S}$, $\mathbf{b}$, and $\boldsymbol\kappa$ when fitted) with
    hold-normalized weights; no prior, regularization or SPD repair enters
    it. The hold-level error is
    $\sigma_\text{obs}^2 = \overline{d_h^2/4} + (3\,\text{mm/s}^2)^2$, where
    $d_h$ is the difference of the two half-hold mean residuals. Its first
    term is measured, and the floor stands for effects a static hold cannot
    reveal. Parameter sigmas are $\sigma_\text{obs}$ times the square roots of
    the diagonal of the inverse information.

    A result is accepted only if the existing coverage/planarity check
    passes on the fit blocks. It also needs the wizard's PolarSPD limits
    (diagonal in $[0.80, 1.25]$, $\kappa(\mathbf{S})\le 6$, off-diagonal RMS
    $\le 0.10$) and every bias below 1.5\,m/s$^2$. The bias sigmas must be
    at most 15\,mm/s$^2$ and the cross-term sigmas at most 0.004. The RMS
    of hold-mean norm errors must be at most 50\,mm/s$^2$, and the held-out
    checks below must pass. No gravity rescaling follows the fit. A
    preliminary static fit after the ten holds runs the same geometry,
    plausibility and held-out checks. If a bias or cross term is
    undetermined, the procedure asks for the main pose that most reduces
    that parameter's variance (rank-one update). If a hold disagrees (its
    studentized hold-out error is over the limit), it recaptures that hold.
    If only the raw hold-out error is over its limit, the other holds
    predict that hold poorly (high leverage), so the procedure adds the pose
    that most reduces the variance of the weakest term instead of repeating
    one. At most two such holds are added; after that the
    calibration fails with the deficient direction named. A diagonal-only
    fallback is never used.

    \subsection{Thermal Bias}
    \label{sec:thermal}

    Every retained block carries the mean of its valid temperatures.
    $T_\text{ref}$ is the mean of the hold-mean temperatures (rounded to
    0.1\,$^\circ$C). The three rechecks follow the gyroscope and
    magnetometer stages in the full wizard, so their temperatures differ by
    whatever the device warmed in normal use. In accelerometer-only mode
    they follow immediately and no wait is inserted. The gyroscope hold is
    not reused as a recheck: the unchanged gyro stage ends after 220
    samples, short of the 5\,s accelerometer hold. The slope is fitted
    jointly with the nine static parameters, from the rechecks' measured
    vectors, never by subtracting repeated poses. It is kept
    (\texttt{LEARNED}) only if every fit block has a valid temperature, the
    hold temperatures span at least 2\,$^\circ$C, each slope sigma after the
    static parameters (the thermal block of the inverse information, i.e.\
    the inverse Schur complement) is at most 1.5\,mm/s$^2$/K, and each
    slope is at most 20\,mm/s$^2$/K. Its runtime clamp is the evidence range
    widened by 5\,$^\circ$C. A single later direction or a span alone does
    not qualify (tests). Otherwise a compatible earlier slope is kept
    (\texttt{PRESERVED}): it must have been learned or preserved by this
    fitter, with bound metadata, for the same ESP32 MAC and IMU type. Only
    the bias is refitted then, so a power-cycle offset is never read as a
    temperature effect. Failing that, $\mathbf{k}=\mathbf{0}$
    (\texttt{UNLEARNED}). Blocks without a temperature are used at
    $T_\text{ref}$. The fit never uses vessel motion or an estimator's
    online bias.

    The model has no thermal scale term. Because the rechecks are
    one-sided, a real thermal scale change aliases into $\mathbf{k}$ along
    the rechecked directions. The campaign reports this case.

    \subsection{Held-Out Checks}

    Leave-one-hold-out: each hold is predicted by a refit without it (frozen
    robust weights), scored only when that refit still determines the model
    (sigmas within three times the gates). The studentized error
    $|e_h| / \sqrt{\max(\sigma_\text{obs}, \sigma_h)^2 +
    \sigma_\text{obs}^2 \mathbf{f}_h^\top \mathbf{C}_{-h}\mathbf{f}_h}$ must
    be at most 6 and $|e_h|$ at most 0.15\,m/s$^2$, where $\sigma_h$ is the
    hold's own split-half scatter and $\mathbf{f}_h$ its regressor. Tilted
    holds carry unique information, so raw hold-out errors mostly measure
    leverage; the studentized form measures consistency. The last 2\,s of
    each recheck never enter the fit or the thermal decision. Their mean
    calibrated norm error, computed with the stored float coefficients, must
    be at most 0.08\,m/s$^2$. It measures repeatability at the recheck, not
    placement independence. On device only norm consistency and model-based
    sigmas are available; three-axis bias accuracy is known only in
    simulation.

    \subsection{Persistence and Runtime}

    The blob (version 3, key \texttt{blob\_m5v3}) holds the accelerometer,
    gyroscope and magnetometer coefficients, the accelerometer runtime
    clamp, qualification metadata and the sensor identity. Only this layout
    and key are read; calibrations saved by earlier firmware are not loaded,
    so such a device runs the wizard again. \texttt{accel\_coeff\_crc} binds the metadata to the
    coefficient set it describes, including \texttt{accel\_g}, the gravity
    the scale was fitted against. The runtime applies an accelerometer
    calibration only when \texttt{accel\_g} equals the firmware's
    $g_\text{cal}$ (to $10^{-5}$\,m/s$^2$); a set fitted against another
    gravity, for example 9.80665\,m/s$^2$, is reported at boot and not
    applied, never rescaled or reinterpreted, while its gyroscope and
    magnetometer calibrations stay in use. The float coefficients are validated after
    the double-to-float conversion, after serialization, and after the
    read-back: finite, exactly symmetric, SPD, within the gates, and
    reproducing the double hold RMS within 0.2\,mm/s$^2$. The save succeeds
    only if the read-back validates and equals the candidate byte for byte,
    so an older blob can never pass as the new one. Accelerometer-only recalibration
    copies the saved gyroscope and magnetometer fields unchanged. On any
    failure or abort nothing is written. \texttt{RuntimeCals} applies the
    calibration once, upstream of every estimator, after
    \texttt{clearM5UnifiedImuCalibration()}.

    \subsection{Validation Campaign}

    \texttt{tests/imu\_calibrate/accel\_cal-test} simulates complete
    sessions for 30 seeds per scenario. It draws a sensor truth: bias
    $\sigma=0.15$\,m/s$^2$, scale 1\,\%, symmetric cross 0.4\,\%, slopes
    2\,mm/s$^2$/K, white and correlated noise, gyro bias. It adds a user who
    reacts, moves with minimum-jerk rotations and places with 2--15$^\circ$
    errors, sometimes first tilting only one axis. A hand-held hold wobbles
    by 0.2--0.6$^\circ$ per mode at 0.2--1.2\,Hz, with drift, 8--12\,Hz
    tremor, 0.5--2\,mm sway and occasional regrips. It also models table
    knocks, stale samples, dropouts and timestamp jitter, and temperature
    histories. The simulated site has the Fair Lawn gravity and the new
    procedure is configured with it, as the wizard is; results are scored
    against the physical truth (a correct calibration maps the site's
    specific force to itself). Scenario \texttt{local\_g} is a
    9.7803\,m/s$^2$ site with a firmware built for it; \texttt{g\_mismatch}
    is the same site with the default Fair Lawn value, and shows the
    expected common scale $g_\text{cal}/9.7803 = 1.00228$ with the bias
    unaffected. The deployed procedure is reproduced: six faces, 6.5\,s wait,
    the first 60 gated samples, \texttt{AccelCalibrator<float,400,1>}
    PolarSPD with its gates and its 9.80665\,m/s$^2$, and the pre-fix
    identity acceptance. The new blocks are also fed to the old fitter
    (with $g_\text{cal}$), and the new procedure runs through the
    same \texttt{AccelCalIo} interface as the wizard. Every result is scored
    through the float runtime on 400 unseen orientations against the truth
    (Table~\ref{tab:campaign}).

    The deployed fit is ill-posed for six clustered poses: its ten-parameter
    quadric takes the cross terms, and through them the centre, from
    in-cluster noise. When its gates rejected the matrix the deployed code
    saved $\mathbf{S}=\mathbf{I}$, $\mathbf{b}=\mathbf{0}$ and reported
    success (column ident.). The new capture alone removes most of the bias
    error; the new fitter adds the determined cross terms, the thermal
    slope, hold balancing and the checks. A design-time sweep of the useful
    hold length (3, 4, 5, 6.5, 8, 10\,s) found lower success at 3--4\,s and
    little accuracy gain beyond 5\,s for 20--40\,s of extra user time, so
    5\,s is used.

    \begin{table*}[t]
        \centering
        \caption{Simulated accelerometer calibration (30 sessions per
        scenario; regenerated by CI from
        \texttt{calibrate\_accel\_campaign\_summary.csv}). Errors at the
        mean capture temperature on unseen orientations through the float
        runtime; cross is the largest symmetric cross-term error; time is
        the user time of the accelerometer holds.}
        \label{tab:campaign}
        \footnotesize
        \IfFileExists{calibrate_accel_campaign_table.pgf}{\input{calibrate_accel_campaign_table.pgf}}{(table generated by the imu\_calibrate CI job)}
    \end{table*}

    \subsection{Limitations}

    All accuracy figures come from simulation; no hardware session or
    recorded log has been evaluated. The raw-sample replay
    (\texttt{accel\_cal-replay}, firmware built with
    \texttt{ATOMS3R\_ICAL\_RAW\_LOG=1}) and the block replay are the route
    for real recordings. The hand, placement and temperature models, the
    documented $x/y$ edge mapping, the sample rate, and the fit time and
    stack use on the ESP32-S3 are unverified on hardware. A learned slope
    describes only its narrow evidence range (typically a few kelvin), not
    full-range thermal behaviour. Thermal scale changes are not modelled.

% -----------------------------------------------------------------------
    \section{Magnetometer Calibration}
    \label{sec:mag_fit}

    \subsection{Numeric Scaling}

    Raw magnetometer values are typically in the range
    $\sim\SI{25}{}$--$\SI{65}{\micro\tesla}$, which leads to condition
    numbers of order $10^6$ in the Gram matrix $\mathbf{H}$
    of \eqref{eq:lsq}.
    The library therefore normalises all samples by a robust radius estimate:
    \begin{equation}
        \mathbf{x}' = \mathbf{x} / \hat{B},
        \quad
        \hat{B} = \mathrm{med}\bigl(\|\mathbf{x}_i - \boldsymbol{\mu}\|\bigr),
    \end{equation}
    fits the unit-sphere model ($R'=1$) in the normalised frame, then
    unscales the result:
    \begin{align}
        \mathbf{b}_{\si{\micro\tesla}} &= \hat{B}\,\mathbf{b}', \\
        \mathbf{A}_{\si{\micro\tesla}} &= \frac{\mathbf{A}'}{{\hat{B}}}\,B_\text{med},
    \end{align}
    where $B_\text{med}$ is the median corrected radius in physical units.

    \subsection{Field Plausibility Check}

    The recovered field magnitude $B_\text{med}$ is accepted only if
    $12 \le B_\text{med} \le \SI{120}{\micro\tesla}$.
    This rejects degenerate solutions and catches situations where
    the sensor was near a strong magnetic disturbance during the
    entire collection.

% -----------------------------------------------------------------------
    \section{Gyroscope Calibration}
    \label{sec:gyro_fit}

    The gyroscope calibrator collects stationary samples gated by:
    \begin{align}
        \|\boldsymbol{\omega}_\text{raw}\| &< \SI{0.07}{\radian\per\second}, \\
        \bigl|\|\mathbf{a}_\text{raw}\| - g\bigr| &< \SI{0.25}{\meter\per\second\squared}.
    \end{align}

    Samples are binned by temperature (same $K=8$ scheme as the accelerometer).
    Within each bin, the mean reading is taken as the bias estimate.
    If at least two bins are populated, a linear temperature model
    \eqref{eq:tempbias} is fitted by ordinary least squares.
    Otherwise a constant bias ($\mathbf{k}=\mathbf{0}$) is used.

% -----------------------------------------------------------------------
    \section{Software Architecture}
    \label{sec:arch}

    \subsection{Namespace and Configuration}

    All symbols are enclosed in namespace \texttt{imu\_cal}.
    The compile-time constant
    \begin{lstlisting}
static constexpr int IMU_CAL_MAX_SAMPLES = 400;
    \end{lstlisting}
    limits the size of all fixed-capacity arrays to prevent dynamic
    allocation.

    \subsection{Error Reporting}

    Failures are reported through a typed enum:
    \begin{lstlisting}
enum class FitFail : uint8_t {
  OK = 0,
  BAD_ARG, TOO_FEW_SAMPLES,
  NON_FINITE_INPUT,
  DEGENERATE_COVERAGE,
  SOLVE_LDLT_FAIL,
  MODEL_LU_NOT_INVERTIBLE,
  MODEL_S_NONPOSITIVE,
  MODEL_SPD_PROJECT_FAIL,
  MODEL_CHOLESKY_FAIL,
  INLIERS_TOO_FEW,
  MAG_FIELD_OUT_OF_RANGE,
  TEMP_BINS_EMPTY,
  ACCEL_S_UNPHYSICAL,
  ACCEL_INFO_LOW,
  ACCEL_CROSSCHECK_FAIL,
  ACCEL_FLOAT_CHECK_FAIL,
};
    \end{lstlisting}
    The helper \texttt{fitFailStr()} converts any code to a printable
    C string for serial diagnostics.

    \subsection{Sample Buffer}

    \texttt{SampleBuffer3<T,N>} is a fixed-capacity ring buffer for
    3-D vectors with associated temperature readings:
    \begin{lstlisting}
template <typename T, int N>
struct SampleBuffer3 {
  Eigen::Matrix<T,3,1> v[N];
  T    tempC[N];
  int  n = 0;
  bool push(const Vec3& x, T tC = T(0));
};
    \end{lstlisting}
    \texttt{push} validates finiteness and rejects samples beyond capacity.

    \subsection{Calibration Result Structs}

    Each sensor has a corresponding result struct:
    \begin{itemize}
        \item \texttt{AccelCalibration<T>}: stores $g$, $\mathbf{S}$,
        \texttt{TempBias3}, RMS, and an \texttt{apply()} method.
        \item \texttt{MagCalibration<T>}: stores $\mathbf{A}$, $\mathbf{b}$,
        $B_\text{field}$, RMS, and an \texttt{apply()} method.
        \item \texttt{GyroCalibration<T>}: stores $\mathbf{S}$,
        \texttt{TempBias3}, and an \texttt{apply()} method.
    \end{itemize}
    All structs carry a Boolean \texttt{ok} flag and use
    \texttt{EIGEN\_MAKE\_ALIGNED\_OPERATOR\_NEW} for correct Eigen alignment.

    \subsection{Calibrator Templates}

    Three template calibrators own a sample buffer and expose a
    two-method API: \texttt{addSample} and \texttt{fit}.

    \begin{lstlisting}
template <typename T, int N, int K_TBINS = 8>
struct AccelCalibrator {
  bool addSample(const Vec3& a, const Vec3& w, T tempC);
  bool fit(AccelCalibration<T>& out, ...);
};

template <typename T, int N>
struct MagCalibrator {
  bool addSample(const Vec3& m_uT);
  bool fit(MagCalibration<T>& out, ...);
};

template <typename T, int N, int K_TBINS = 8>
struct GyroCalibrator {
  bool addSample(const Vec3& w, const Vec3& a, T tempC);
  bool fit(GyroCalibration<T>& out, ...);
};
    \end{lstlisting}

% -----------------------------------------------------------------------
    \section{Typical Usage}
    \label{sec:usage}

    The following snippet illustrates a typical calibration workflow
    for a float-precision magnetometer with 200 sample capacity:

    \begin{lstlisting}
#include "imu_cal.h"
using namespace imu_cal;

MagCalibrator<float, 200> cal;

// -- collection phase (rotate device freely) --
while (!cal.buf.n < 100) {
  Eigen::Vector3f m_raw;
  readMag(m_raw);          // application-specific
  cal.addSample(m_raw);
}

// -- fitting phase --
MagCalibration<float> result;
FitFail why;
if (!cal.fit(result, 3, 0.15f, 1e-6f, &why)) {
  Serial.println(fitFailStr(why));
  return;
}

// -- application phase --
Eigen::Vector3f m_cal = result.apply(m_raw);
    \end{lstlisting}

    The accelerometer and gyroscope calibrators additionally require a
    simultaneous gyroscope / accelerometer reading for their respective
    static-motion gates.

% -----------------------------------------------------------------------
    \section{Numerical Notes}
    \label{sec:numerical}

    \subsection{Precision}

    The library is parameterised by scalar type \texttt{T}.
    \texttt{float} (32-bit) is sufficient for field deployment.
    \texttt{double} (64-bit) is recommended when validating the algorithm
    on a host machine.
    Internally, all \texttt{sqrt}, \texttt{fabs}, and \texttt{isfinite}
    calls are cast through \texttt{double} to avoid platform-specific
    float precision issues on AVR-class targets.

    \subsection{No Dynamic Allocation}

    All internal arrays (residuals, inlier mask, sorted quantile buffer)
    are declared at fixed size \texttt{IMU\_CAL\_MAX\_SAMPLES}.
    The full-matrix accelerometer fitter keeps its per-block workspace in the
    fitter object and solves one fixed-size $12\times12$ LDLT; parameters
    that are not free are pinned by identity rows.
    The AccelCalibrator keeps its per-bin index arrays in a mutable
    \texttt{Workspace} member to reduce stack pressure inside \texttt{fit()}.

    \subsection{Eigen Dependency}

    The header conditionally includes either \texttt{<Eigen/Dense>} and
    \texttt{<Eigen/Eigenvalues>} (host builds) or
    \texttt{<ArduinoEigenDense.h>} (Arduino builds):
    \begin{lstlisting}
#ifdef EIGEN_NON_ARDUINO
#  include <Eigen/Dense>
#  include <Eigen/Eigenvalues>
#else
#  include <ArduinoEigenDense.h>
#endif
    \end{lstlisting}
    No other external dependencies are required.

% -----------------------------------------------------------------------
    \section{Failure Mode Summary}

    Table~\ref{tab:failures} summarises the failure codes, their root
    causes, and typical remedies.

    \begin{table}[t]
        \centering
        \caption{Failure codes and remedies}
        \label{tab:failures}
        \renewcommand{\arraystretch}{1.2}
        \begin{tabular}{@{}lp{4.2cm}@{}}
            \toprule
            \textbf{Code} & \textbf{Meaning / Remedy} \\
            \midrule
            \texttt{TOO\_FEW\_SAMPLES}      & Collect at least 80 samples per bin \\
            \texttt{NON\_FINITE\_INPUT}     & Check ADC / I$^2$C link for errors \\
            \texttt{DEGENERATE\_COVERAGE}   & Rotate more; cover all six hemisphere faces \\
            \texttt{SOLVE\_LDLT\_FAIL}      & Increase ridge; unlikely with normal data \\
            \texttt{MODEL\_S\_NONPOSITIVE}  & Sensor may be saturated; check range \\
            \texttt{MODEL\_SPD\_PROJECT\_FAIL} & Near-degenerate geometry; collect more data \\
            \texttt{MODEL\_CHOLESKY\_FAIL}  & SPD projection failed; rarely occurs \\
            \texttt{INLIERS\_TOO\_FEW}      & Too many outliers; check mounting, vibration \\
            \texttt{MAG\_FIELD\_OUT\_OF\_RANGE} & Keep away from strong magnets; recheck sensor \\
            \texttt{TEMP\_BINS\_EMPTY}      & Too few samples per temperature bin; slow down cooling \\
            \texttt{ACCEL\_S\_UNPHYSICAL}   & Diagonal or condition check failed; check gates \\
            \texttt{ACCEL\_INFO\_LOW}      & Holds do not determine a bias or cross term; extra hold \\
            \texttt{ACCEL\_CROSSCHECK}     & A held-out hold disagrees; recapture it (raw-only: extra hold) \\
            \texttt{ACCEL\_FLOAT\_CHECK}   & Stored float coefficients failed re-validation \\
            \bottomrule
        \end{tabular}
    \end{table}

% -----------------------------------------------------------------------
    \section{Conclusion}

    \texttt{imu\_cal} provides in-situ IMU calibration in C\texttt{++}/Eigen
    headers. Magnetometer and library accelerometer calibration use a
    robust 10-parameter ellipsoid fit; gyroscope calibration fits a linear
    temperature-bias model from stationary data. The AtomS3R accelerometer
    procedure uses guided static holds chosen for observability, a geometric
    full-matrix fit with a qualified thermal slope, held-out checks,
    re-validated float coefficients and verified persistence.

% -----------------------------------------------------------------------
    \begin{thebibliography}{9}

        \bibitem{Frosio2009}
        I.~Frosio, F.~Pedersini, and N.~A.~Borghese,
        ``Autocalibration of MEMS Accelerometers,''
        \emph{IEEE Trans.\ Instrumentation and Measurement},
        vol.~58, no.~6, pp.~2034--2041, Jun.\ 2009.

        \bibitem{Vasconcelos2011}
        J.~F.~Vasconcelos, G.~Elkaim, C.~Silvestre, P.~Oliveira, and B.~Cardeira,
        ``Geometric Approach to Strapdown Magnetometer Calibration in Sensor Frame,''
        \emph{IEEE Trans.\ Aerospace and Electronic Systems},
        vol.~47, no.~2, pp.~1293--1306, Apr.\ 2011.

        \bibitem{Eigen}
        G.~Guennebaud, B.~Jacob, et al.,
        ``Eigen v3,''
        \url{http://eigen.tuxfamily.org}, 2010.

        \bibitem{Higham1988}
        N.~J.~Higham,
        ``Computing a Nearest Symmetric Positive Semidefinite Matrix,''
        \emph{Linear Algebra and its Applications},
        vol.~103, pp.~103--118, 1988.

    \end{thebibliography}

\end{document}

